% =============================================================================
% §1 Introduction — phenomenon (R1); puzzle; contributions C1–C5; the arc;
% the roadmap / division-of-labor paragraph (Fina §3, mandate A).
% =============================================================================
\section{Introduction}
\label{sec:intro}

Between 2015 and 2024, several Latin American countries experienced fertility
declines of a speed with few peacetime precedents. Costa Rica's total fertility
rate (TFR) fell from 1.83 in 2010 to 1.12 in 2024 --- a 39 percent decline, most
of it after 2016. Colombia's national registry series fell from 1.7 to 1.1 in
under a decade; Chile reached 1.03 in 2024. These are not marginal adjustments
within a below-replacement regime; they are collapses to among the lowest
fertility levels ever recorded, in middle-income countries, within roughly
eight years --- the kind of tempo that has moved global demographic
projection debates from convergence scenarios to genuine uncertainty about
the fertility floor \citep{fernandezvillaverde2026}. Remarkably, the acceleration is barely visible in the smoothed
international series most analysts consult: World Bank and World Population
Prospects figures, which interpolate through model schedules, trail the national
registries by up to half a child per woman at the end of the sample
(Section~\ref{sec:phenomenon}).

The speed is the puzzle. Classical transition frameworks --- including the Second
Demographic Transition (SDT) that correctly anticipated the \emph{direction} of
change \citep{lesthaeghe1986twee, vandekaa1987, lesthaeghe2010} --- describe
ideational and institutional change diffusing across generations. They give no
mechanism for a national TFR losing a third of its level in eight years. Fast
declines invite fast-dynamics explanations, and the current literature and public
debate lean readily on \emph{reflexive} accounts: social-media-amplified norm
cascades, coordination tipping points, self-reinforcing childlessness. These
mechanisms are theoretically attractive precisely because they are fast. Whether
the data actually exhibit their signature is the question this paper answers.

We make five contributions. \textbf{First} (measurement), we document the
collapse against national vital registries and household microdata rather than
smoothed international series, showing that raw annual births fell 18--47 percent
across the collapse countries while the female population of reproductive age
grew --- the decline is behavioral, not a denominator artifact --- and that the
standard international series systematically erase its acceleration.
\textbf{Second} (state variable), we show that the proximate driver is union
\emph{composition}: total union prevalence among women 20--39 is roughly flat
while marriage collapses and cohabitation substitutes --- in Colombia, marriage
halves (20.7 to 11.9 percent) on a flat union plateau. \textbf{Third}
(identification), we discriminate between candidate mechanisms with a
pre-registered, outcome-robust design: age--period--cohort (APC) curvature
localizes the change, and the \emph{sign} of state dependence in pseudo-cohort
panels tests for reflexivity. The verdict is unambiguous: the state-dependence
coefficient is positive --- stabilizing, the opposite of a cascade signature ---
in every specification, in both countries; the structure is \emph{cohort
replacement}, with a small exogenous period component in Colombia.
\textbf{Fourth} (sufficiency), we show that a deliberately transparent cohort
microsimulation with four interpretable parameters and no social feedback is
\emph{sufficient} for the observed composition collapse --- magnitude and
late-loading --- and that its single calibrated cohort inflection point
independently recovers the APC crossover it was never fit to. \textbf{Fifth}
(interpretation), we frame the findings as \emph{channel-assisted compression}
of the SDT: change concentrates where the institutional channel --- Latin
America's long-standing consensual-union system \citep{esteve2012} --- was
already wide, which yields a falsifiable out-of-sample prediction for
marriage-dominant regimes.

The paper's arc is deliberately outcome-robust. We built the fashionable model
first: an agent-based simulation with a social-norm coordination threshold, the
formalization of the cascade narrative. The data rejected it twice --- the
calibrated threshold produced a smooth glide rather than the observed
acceleration, and the econometric cascade test returned the wrong sign in every
specification. What survived is simpler: successive cohorts arrive with lower
marriage propensity, and the stock composition follows as they flow through the
age distribution. We report the rejections as prominently as the survivor,
because the discrimination is what makes the survivor credible.

A roadmap, which is also the paper's division of labor. Registries and microdata
\emph{measure} the phenomenon (Section~\ref{sec:phenomenon}).
Section~\ref{sec:mechanisms} states the three candidate mechanisms and their
observable signatures. APC curvature and state dependence \emph{identify} the
mechanism econometrically (Section~\ref{sec:discrimination}); identification is
econometric throughout, and the simulation model is never asked to perform it.
The agent-based model \emph{tests sufficiency} of the identified mechanism ---
whether cohort replacement alone can generate what we observe --- under a
pre-registered evaluation (Section~\ref{sec:model}). The SDT \emph{frames} the
result and generates the out-of-sample prediction
(Section~\ref{sec:interpretation}). Section~\ref{sec:limitations} states the
scope conditions and data limitations; Section~\ref{sec:conclusion} concludes.
